\documentclass[a4paper,11pt]{article}
\usepackage[utf8]{inputenc}
\usepackage[T1]{fontenc}
\usepackage[ngerman]{babel}
\usepackage{amsmath,amssymb}
\usepackage[margin=2.3cm]{geometry}
\usepackage{enumitem}
\usepackage{fancyhdr}
\usepackage{array}
\usepackage[table]{xcolor}
\usepackage{tikz}
\definecolor{bgdark}{HTML}{1E1E1E}
\definecolor{fgtext}{HTML}{E6E6E6}
\definecolor{gridcol}{HTML}{4A4A4A}
\pagecolor{bgdark}
\arrayrulecolor{fgtext}
\AtBeginDocument{\color{fgtext}}
\renewcommand{\headrule}{\color{fgtext}\hrule height \headrulewidth}
\newcommand{\karo}[1]{\par\noindent
  \begin{tikzpicture}
    \draw[step=5mm, gridcol, very thin] (0,0) grid (\linewidth,#1);
  \end{tikzpicture}\par}

\newcommand{\diff}[1]{\hfill\normalfont\textbf{[#1]}}
\newcommand{\aufg}[2]{\subsection*{Aufgabe #1 \diff{#2}}}
\renewcommand{\arraystretch}{1.2}

\pagestyle{fancy}
\fancyhf{}
\lhead{RADM -- Logik}
\rhead{Übungsblatt 01}
\cfoot{\thepage}

\begin{document}

\begin{center}
  {\LARGE\bfseries Übungsaufgaben Logik}\\[2pt]
  {\large Relationen, Algebra, Diskrete Mathematik -- Teil 2}\\[10pt]
\end{center}

\noindent
\begin{tabular}{@{}p{0.6\textwidth}@{}p{0.4\textwidth}@{}}
  \textbf{Name:} \hrulefill & \textbf{Datum:} \hrulefill
\end{tabular}

\vspace{4pt}
\noindent\rule{\textwidth}{0.4pt}
\vspace{4pt}

\noindent\textit{Themen: Aussagen, Wahrheitstafeln, Implikation/Äquivalenz, Tautologie,
logische Gesetze, Normalformen (DN/KN/kDN), Resolution, Prädikatenlogik.\\
Schwierigkeit: \textbf{[leicht] / [mittel] / [schwer]}. Notation wie in der Vorlesung
(Wahrheitswerte $w$/$f$).}

\vspace{6pt}

% ---------------------------------------------------------------
\aufg{1}{leicht}
Welche der folgenden Sätze sind \emph{Aussagen} im Sinne der Aussagenlogik?
Begründe jeweils kurz.
\begin{enumerate}[label=(\alph*)]
  \item Wien ist die Hauptstadt von Österreich.
  \item $x + 3 = 7$.
  \item Wie spät ist es?
  \item $5$ ist eine Primzahl.
  \item Dieser Satz besteht aus fünf Wörtern.
  \item Bitte schließe die Tür.
\end{enumerate}
\karo{4.5cm}

% ---------------------------------------------------------------
\aufg{2}{leicht}
Erstelle die vollständige Wahrheitstafel für
\[ F(a,b) = (a \wedge \neg b) \vee (\neg a \wedge b). \]
Welche bekannte Verknüpfung beschreibt $F$?
\karo{6cm}

% ---------------------------------------------------------------
\aufg{3}{mittel}
Zeige mit Hilfe von Wahrheitstafeln:
\begin{enumerate}[label=(\alph*)]
  \item $a \to b$ ist äquivalent zur Kontraposition $\neg b \to \neg a$.
  \item $a \to b$ ist \emph{nicht} äquivalent zur Umkehrung $b \to a$.
\end{enumerate}
\karo{7cm}

% ---------------------------------------------------------------
\aufg{4}{mittel}
Klassifiziere jede Formel als \emph{Tautologie}, \emph{widerspruchsvoll}
oder (nur) \emph{erfüllbar}. Begründe (z.\,B. mit Wahrheitstafel).
\begin{enumerate}[label=(\alph*)]
  \item $(a \to b) \leftrightarrow (\neg a \vee b)$
  \item $a \wedge \neg a$
  \item $(a \vee b) \wedge \neg a$
  \item $\big((a \to b) \wedge a\big) \to b$
\end{enumerate}
\karo{6.5cm}

% ---------------------------------------------------------------
\aufg{5}{mittel}
Vereinfache den folgenden Term mit den logischen Gesetzmäßigkeiten
(gib bei jedem Schritt das verwendete Gesetz an):
\[ F = \neg(\neg a \vee b) \vee (a \wedge b). \]
\karo{6cm}

% ---------------------------------------------------------------
\aufg{6}{mittel}
Gegeben sei $F(a,b,c) = (a \to b) \wedge c$.
\begin{enumerate}[label=(\alph*)]
  \item Gib eine disjunktive Normalform (DN) von $F$ an.
  \item Gib eine konjunktive Normalform (KN) von $F$ an.
\end{enumerate}
\karo{6cm}

% ---------------------------------------------------------------
\aufg{7}{schwer}
Bestimme die \emph{kanonische disjunktive Normalform} (kDN) von
\[ F(a,b,c) = a \vee (b \wedge \neg c). \]
\karo{8cm}

% ---------------------------------------------------------------
\aufg{8}{schwer}
Zeige mittels \emph{Resolution}, dass die Schlussfolgerung
\[ \{\, a \to b,\; b \to c,\; a \,\} \models c \]
korrekt ist. Gib alle Resolutionsschritte an.
\karo{7cm}

% ---------------------------------------------------------------
\aufg{9}{schwer}
Ein vereinfachter ,,Mini-Kriminalfall``. Aus den Zeugenaussagen ergeben sich die
Klauseln
\[ \{\, \neg a \vee b,\;\; \neg b \vee c,\;\; a \vee d,\;\; \neg d \,\}. \]
Zeige mittels Resolution, dass daraus $c$ folgt, d.\,h.
$\{\neg a \vee b,\ \neg b \vee c,\ a \vee d,\ \neg d\} \models c$.
\karo{7.5cm}

% ---------------------------------------------------------------
\aufg{10}{mittel}
Formalisiere die folgenden Aussagen prädikatenlogisch mit Quantoren.
Verwende die angegebenen Prädikate.
\begin{enumerate}[label=(\alph*)]
  \item ,,Alle Studierenden bestehen die Klausur.`` \quad ($S(x)$: $x$ ist Studierende(r); $B(x)$: $x$ besteht die Klausur)
  \item ,,Es gibt eine natürliche Zahl, die gerade und prim ist.``
  \item ,,Kein Mensch kann fliegen.`` \quad ($M(x)$: $x$ ist ein Mensch; $F(x)$: $x$ kann fliegen)
  \item ,,Zu jeder reellen Zahl gibt es eine größere.``
\end{enumerate}
\karo{6cm}

\end{document}
